% =====================================================
% Lời nói đầu (báo cáo học phần, rút gọn)
% =====================================================
\UnnumberedChapter{LỜI NÓI ĐẦU}

Tin tức giả mạo hiện nay ngày càng lan tràn như là hệ lụy tất yếu của thời đại số. Người đọc khi tiếp cận thông tin trên mạng xã hội, các nền tảng báo điện tử có thể gặp phải những bài viết xuyên tạc, thổi phồng, đánh lạc hướng dư luận, gây hoang mang xã hội hoặc trục lợi thương mại. Việc kiểm chứng thủ công một bài tin tốn rất nhiều thời gian, trong khi số lượng bài viết được đăng tải mỗi ngày là vô cùng lớn. Nếu không có công cụ hỗ trợ một cách nhanh chóng, chính xác và hiệu quả thì người đọc rất khó phân biệt được tin thật và tin giả.

Xuất phát từ thực tế đó, nhóm 4 sinh viên lớp 18-11 đã lựa chọn đề tài: ``Phát hiện tin tức giả mạo sử dụng thuật toán KNN''. Mô hình áp dụng thuật toán láng giềng gần nhất (KNN --- K-Nearest Neighbors) trên đặc trưng TF-IDF với bộ dữ liệu ISOT Fake News gồm hơn 44.000 bài viết; kết quả đánh giá trên tập kiểm tra đạt độ chính xác 89,35\%, đồng thời được so sánh với các mô hình học máy truyền thống khác (Naive Bayes, Logistic Regression, Passive Aggressive) trên cùng dữ liệu. Nhóm cũng tiến hành chuỗi thí nghiệm tinh chỉnh các siêu tham số của KNN (số láng giềng k, phép đo khoảng cách, trọng số láng giềng) và đóng gói giải pháp thành ứng dụng demo cho phép người dùng kiểm tra nhanh một bài tin tức bất kỳ.

Nhóm xin chân thành cảm ơn \GiangVienHD{} đã tận tình hướng dẫn và hỗ trợ trong suốt quá trình thực hiện đề tài. Do thời gian và kiến thức còn hạn chế nên báo cáo không tránh khỏi thiếu sót, nhóm rất mong nhận được góp ý của cô để hoàn thiện hơn.

\begin{flushright}
Nhóm xin chân thành cảm ơn!
\end{flushright}
